\section{The eleven fields of the CM literature input}

% BEGIN C72M1
\paragraph{C72M1. Good reduction and inertia.}
Let $R$ be the given discrete valuation ring, $K$ its fraction field,
and $I$ the inertia group at the chosen prolongation of its valuation to
a separable closure. For $m$ nonzero in the residue field,
the geometric $m$-torsion is precisely the kernel of $[m]$ on
the separable-closure points. Passing to the inverse limit for
$m=\ell^n$ gives the Tate module with the indicated Galois action.
Serre and Tate, \emph{Good reduction of abelian varieties},
Theorem~1, give all four equivalences over a perfect residue field:
an abelian model over $R$, trivial inertia on every such torsion group,
trivial inertia for infinitely many such $m$, and trivial inertia
on one prime-to-residue-characteristic Tate module. Theorem~2(i)
identifies potential good reduction with finite inertia image on that
module. The good model is taken at a prolongation to a finite extension,
as required by the quantifiers on its valuation ring.

For the CM clause, the given unital action of a ring $O$ with fraction
field $F$ induces $F\to\operatorname{End}^0_K(A)$. Indeed
$O\otimes\mathbb Q=F$: it is a rational subalgebra of a number
field, hence a finite-dimensional domain and therefore a field,
and it contains $O$. The action extends by rational linearity.
The prescribed degree $[F:\mathbb Q]=2\dim A$ now gives exactly the
CM hypothesis of Serre--Tate, Theorem~6(a), which gives potential good
reduction over the given finite residue field. This supplies
\texttt{st} in all three parts.
% END C72M1

% BEGIN C72M2
\paragraph{C72M2. The finite idèle in the CM torsion diagram.}
Embed the given imaginary quadratic field $K$ in $\mathbb C$ and
identify the analytic elliptic curve $X$ with
$\mathbb C/\mathfrak a$ using the specified $\xi$. Its stipulated
analyticity off the origin extends across the origin by translation
and the group law. A bijective holomorphic homomorphism of these
complex Lie groups is an isomorphism. The torsion is
$K/\mathfrak a$, canonically identified with
$\mathbb A_{K,f}/\widehat{\mathfrak a}$ by approximation.
The CM type and its reflex type for this chosen embedding each have
one element, and the reflex norm is the identity.
Apply the fundamental theorem of complex multiplication over the
reflex field (Shimura--Taniyama, or Milne,
\emph{The fundamental theorem of complex multiplication}, \S3).
The source sends a prime uniformizer to arithmetic Frobenius.
In this convention, if $\sigma|_{K^{\rm ab}}=\operatorname{Art}(s)$,
the conjugate lattice is $s_f^{-1}\mathfrak a$ and the map on
torsion in the two uniformizations is multiplication by $s_f^{-1}$
in the finite adelic quotients. The inverse at an unramified prime
agrees with the isogeny from the lattice $\mathfrak a$ to the
larger lattice $\mathfrak p^{-1}\mathfrak a$. The theorem also
specifies the unit components on torsion. It applies to any supplied
Artin homomorphism since its finite abelian restrictions agree with
arithmetic reciprocity on prime ideals away from the conductors.
Indeed take a common ray modulus for the two continuous finite
quotients; ideals away from it generate the ray group. The choice of embedding used to
describe restriction does not change the abelianized action. The
resulting $\xi'$ is unique because torsion is dense and holomorphic
uniformizations are continuous. This proves \texttt{cm}.
% END C72M2

% BEGIN C72M3
\paragraph{C72M3. Singular moduli and the Hilbert class field.}
For two fractional ideals the associated $j$-invariants coincide
exactly when the complex elliptic curves are isomorphic, equivalently
when the lattices are homothetic. If $z\mathfrak a=\mathfrak b$,
taking any nonzero $a\in\mathfrak a$ gives $z=(za)/a\in K^\times$.
Thus the distinct values $j(\mathfrak a)$ are indexed by the ideal
class group, including the exceptional lattices with extra units.
By C72M2 the action on these values is inverse ideal translation.
It is transitive and its stabilizer for $j(\mathcal O_K)$ is exactly
the preimage of the principal class. Class field theory identifies
the corresponding fixed field with the Hilbert class field:
all finite local units are killed and the quotient is the ordinary
ideal class group. There are no real places. In particular
$K(j(\mathcal O_K))/K$ has degree $h_K$, is maximal among finite
unramified abelian subextensions of $\mathbb C/K$, and the orbit of
its distinguished generator is precisely the set of $h_K$ ideal
invariants. This gives all five clauses of \texttt{hj}.
% END C72M3

% BEGIN C72M4
\paragraph{C72M4. Rational ray values and the smoothed difference.}
For a nonzero finite modulus $\mathfrak f$ the source uses the
narrow ray class group, with all real places in the modulus.
Indeed its relation identifies $I,J$ when $aI=bJ$ for totally
positive $a,b$ that are both $1$ modulo $\mathfrak f$.
Conversely a ray principal quotient may be written in this form by
clearing its denominators away from $\mathfrak f$ with a totally
positive denominator congruent to $1$ there (ideal approximation).
Expand the given rational function $\varepsilon$ on this finite
group in class indicators. The associated partial zeta series
converge absolutely for $\Re s>1$, continue holomorphically off
$1$, and have rational values at $1-k$ for all $k\ge1$ by
Siegel--Klingen. Their finite linear combination gives the first
half of \texttt{dri}.

For the second half take $\mathfrak f=(\ell^n)$, $n\ge1$. All Euler
factors above $\ell$ are already omitted. On the tower of ray groups,
the prime-to-$\ell$ ideal $\mathfrak c$ has norm character value
$N\mathfrak c$ and translates the test function by
$x\mapsto\varepsilon(\mathfrak c x)$. Deligne--Ribet,
\emph{Values of abelian $L$-functions at negative integers over
totally real fields}, Theorem~(0.4), applied to the single weight
$k$, gives
\[
 L(1-k,\varepsilon)-(N\mathfrak c)^k L(1-k,\varepsilon_{\mathfrak c})
       \ \in\ \mathbb Z_\ell .
\]
Its integrality hypothesis holds since both the test values and the
$\ell$-adic unit norm powers are integral. This includes $\ell=2$.
For the rational field the same one-variable Kummer congruence
appears in \S1 of that paper.
Analytic uniqueness off $1$ identifies the two supplied ray values
with these evaluations. The target norm inequality means membership
in $\mathbb Z_\ell$.
% END C72M4

% BEGIN C72M5
\paragraph{C72M5. The ordinary ideal class paired by Poisson summation.}
Choose a fractional ideal $\mathfrak b$ in class $C^{-1}$.
The nonzero elements $x\in\mathfrak b$, modulo all units, give the
integral ideals $(x)\mathfrak b^{-1}$ in class $C$, with norm
$|N x|/N\mathfrak b$. The trace-dual lattice is
$\mathfrak b^*=\mathfrak D^{-1}\mathfrak b^{-1}$, where
$\mathfrak D$ is the different. Hence the series from the dual
lattice is indexed by the class
$[\mathfrak b^*]^{-1}=[\mathfrak D]C^{-1}$, precisely $C'$.

Here are the factors in Hecke's theta integral. Put
$d=|\operatorname{disc}F|$ and $r=[F:\mathbb Q]$. Embed
$\mathfrak b$ in $\mathbb R^r$ and rescale it by the positive real
scalar $(\sqrt d\,N\mathfrak b)^{-1/r}$, giving a lattice $L$ of
covolume one. Its Euclidean dual has the same covolume-one normalization
of $\mathfrak b^*$. For positive diagonal $t$, Poisson summation for
the Gaussian gives
$\theta_L(t)=(\prod t_i)^{-1}\theta_{L^*}(t^{-1})$, with
$\theta_L(t)=\sum_{x\in L}\exp(-\pi\sum_i t_i^2x_i^2)$.
Integrate $\theta_L(t)-1$ against
$(\prod t_i)^s\prod_i dt_i/t_i$, modulo the absolute values
of the units. Unfolding a nonzero orbit gives one
$\pi^{-s/2}\Gamma(s/2)$ per real place (up to a factor $2^{-r}$
and the common torsion-unit multiplicity). The coordinate rescaling
supplies $d^{s/2}(N\mathfrak b)^s$.
Splitting at $\prod t_i=1$ and applying Poisson summation to the
lower part exchanges $(L,s)$ with $(L^*,1-s)$. The zero terms have
the same poles at $0,1$ on both sides. Thus the two completed series
$d^{s/2}\Gamma_{\mathbb R}(s)^r Z_C(s)$ and its counterpart at
$1-s$ agree meromorphically. Dividing by the gamma factors leaves
$Z_C$ holomorphic off $1$. At all the source's allowed arguments both
gamma products are finite and neither pole contributes. This gives
\texttt{pz}.
% END C72M5

% BEGIN C72M6
\paragraph{C72M6. Hecke's functional equation in the prescribed normalization.}
At the complex places the continuous unitary character has angular
exponents $n_i$. The asserted formula on $E^\times$ fixes the whole
infinite component by weak approximation. The type selects one
embedding per conjugate pair for an arbitrary totally complex $E$.
For the standard additive character, trivial on the diagonal field,
the infinite component is $e^{-2\pi i\operatorname{Tr}(x)}$.
Tate's local and global zeta integral theorem
(\emph{Fourier analysis in number fields and Hecke's zeta-functions})
has the following factors with self-dual additive measures.
At a complex place the Gaussian times the harmonic monomial of degree
$|n_i|$ gives $2(2\pi)^{-s-|n_i|/2}\Gamma(s+|n_i|/2)$.
Fourier transformation with the stated trace character exchanges the
two angular characters and multiplies the harmonic Gaussian by
$(-i)^{|n_i|}$.

At a finite place write $a$ for the conductor exponent of $\chi$,
$d_v$ for the exponent of the different, $q$ for the residue cardinality
and $c=\varpi^{a+d_v}$. The annihilator of $\mathcal O_v$ is
$\mathfrak p^{-d_v}$, so that $\mathcal O_v$ has self-dual volume
$q^{-d_v/2}$. If $a>0$, Fourier-transform
$\chi_v^{-1}\mathbf 1_{\mathcal O_v^\times}$. The integral vanishes
off the shell $c^{-1}\mathcal O_v^\times$ by primitivity, and at
$c^{-1}$ it is
$q^{-a-d_v/2}\sum_{u\in(\mathcal O_v/\mathfrak p^a)^\times}
\chi_v(u)^{-1}\psi_v(u/c)$. Taking its Mellin transform at $1-s$
contributes $\chi_v(c)q^{(a+d_v)(1-s)}$. Thus the local factor is
\[
 \chi_v(c)q^{-s(a+d_v)+d_v/2}
 \sum_{u\in(\mathcal O_v/\mathfrak p^a)^\times}
       \chi_v(u)^{-1}\psi_v(u/c).
\]
The summand is independent of the lift since both characters are
unchanged by the indicated congruence. For $a=0$ Fourier-transform
$\mathbf 1_{\mathcal O_v}$ instead and divide out the two unramified
Euler factors. The same prefactor results, with sum replaced by $1$.
It is $1$ if also $d_v=0$, so the ramified and different primes
account for the entire finite product.

The adelic Poisson formula now gives
$\Lambda_\chi(s)=(-i)^{\sum_i |n_i|}\prod_v\epsilon_v(s)
\Lambda_{\chi^{-1}}(1-s)$. At least one nonzero angular exponent
rules out a power of the adelic norm, so both completed functions
are entire. Where $\Re s>1$ the unramified Euler products converge
absolutely and have the stated complex gamma factors. This supplies
the full totally complex field \texttt{fe}, including odd exponents.
% END C72M6

% BEGIN C72M7
\paragraph{C72M7. Periods, tangent vectors and the integral Chern class.}
Let $u$ be the prescribed holomorphic uniformization and $\Lambda$
its kernel. A full discrete subgroup of $\mathbb C^g$ is free of rank
$2g$, and $u$ identifies $\mathbb C^g/\Lambda$ with $A(\mathbb C)$.
For $x\in\Lambda$ integrate degree-one de Rham classes on the loop
$t\mapsto u(tx)$, $0\le t\le1$. D72P3 and the algebraic de Rham
comparison over $\mathbb C$ identify this functional with an element
$P(x)$ of the dual used in the source. All degree-one classes are
primitive: the Künneth decomposition for the group law has two terms
and restriction to either axis is the identity. Integration identifies
$\Lambda\otimes\mathbb C$ with the full dual, giving the span,
dimension and additivity. A holomorphic one-form is invariant
under translations, so its pullback to the covering vector space
is constant. Its integral along this loop is its evaluation at
$du_0(x)$. This is the specified tangent-vector test: on a pair
of first-order lifts the infinitesimal difference of their form
integrals is evaluation of the form on their difference. Composing
the loop with an endomorphism carrying $x$ to $Mx$ gives the
required covariance.

For the factor of automorphy of a given Hermitian matrix $H$ the
alternating integral form of Appell--Humbert is
$\operatorname{Im} H(x,y)$, up to the fixed order and frame convention.
Indeed, subtract the changes of its logarithm around the two sides of
the parallelogram $(x,y)$. The constant unit factors cancel, leaving
$\pi(H(x,y)-H(y,x))=2\pi i\operatorname{Im}H(x,y)$, up to that
order. The source's divided-power exponential cocycle for the first
Chern class becomes the usual logarithmic coboundary over $\mathbb C$.
Under de Rham comparison it is $2\pi i$ times the integral first
Chern class with the corresponding choice of sign. The alternating
evaluation of each cup $a\smile b$ on $(x,y)$ is
$P(x)(a)P(y)(b)-P(y)(a)P(x)(b)$, exactly the prescribed wedge
formula. The ordering of the coboundary and the frame convention
are fixed once for all groups and bundles. Their combined sign
therefore supplies the single $\epsilon\in\{1,-1\}$ in \texttt{per}.
% END C72M7

% BEGIN C72M8
\paragraph{C72M8. A family of periods and its value on an infinitesimal point.}
On the given open set the evaluation functions define a holomorphic map
to $(\operatorname{Spec} B)^{\rm an}$: one can check this on a finite
set of algebra generators. Pull back the smooth proper family to a
small simply connected ball $V$ around zero. Trivialize its holomorphic
Lie bundle using a frame agreeing at zero with the differential of
the specified $u_0$. The fibrewise exponential gives the required
uniformizations. The kernel is the locally constant lattice of
integral first homology inside the Lie bundle; its sections are
holomorphic by the local inverse theorem for the exponential. Continue
the marked lattice in that local system to obtain $\tau_z$. At zero
the two given models are canonically identified as pullbacks, so the
exponentials agree and $\tau_0$ fixes each lattice element. The
pullback of every $L_z$ to $\mathbb C^g$ is holomorphically trivial
by the exponential sequence and Stein vanishing. A flat family of
these line bundles is unnecessary since their frames are chosen
separately in each fibre. Off $V$ the quantified uniformizations and
lattice isomorphisms exist fibrewise. On $V$ use precisely the
period functional and the common sign of C72M7. A global group
endomorphism acts holomorphically on the Lie bundle and on the
integral local system. Its derivative intertwines the exponentials,
and preserves every transported lattice relation known at zero.

Under the relative de Rham comparison, integration against a continued
cycle is a horizontal functional for the Gauss--Manin connection
(Katz--Oda, \emph{On the differentiation of de Rham cohomology
classes with respect to parameters}). Pairing it with the given
relative class $c$ thus gives a holomorphic function $p_c$ on $V$.
For the last equality take
$J_N=\mathbb C[z_i:i\in I]/(z_i)^{N+1}$, where $(z_i)$ denotes the
ideal generated by all coordinates. The product rule makes truncated
Taylor expansion of $b\mapsto y(z)(b)$ an algebra map $B\to J_N$.
The prescribed tuple defines the
homomorphism $J_N\to D$, $z_i\mapsto t_i$ even when $D$ has no
finiteness hypothesis. It preserves the divided powers since on a
rational algebra $\gamma_r(a)=a^r/r!$. The family over $D$ is pulled back from the
$N$-jet because the equality for $j(b)$ holds on every $b\in B$.
The infinitesimal stratification of the connection transports the
central functional to this jet. In characteristic zero its Taylor
isomorphism is determined by iterating the connection; applied to
$c$ its coefficients are exactly the ordinary derivatives of $p_c$
divided by the factorials. The evaluation on the constant central
fibre is fixed by the predicate of correspondence: on each scalar
extension of a central cohomology class it equals its prescribed
central period. Such classes span by base change on the constant
family, and the crystal comparison identifies their transport to
the deformed fibre. Thus the stipulated $d$ gives this very
horizontal extension. The required result is
\[
 \sum_{k=0}^N\frac1{k!}\!
  \sum_{(i_1,\ldots,i_k)\in I^k}
  (\partial_{i_1}\cdots\partial_{i_k}p_c)(0)\,t_{i_1}\cdots t_{i_k}.
\]
The comparison along the two lifts commutes with pullback along
$J_N\to D$, so evaluating at this tuple proves the formula for
every permitted $D$ in \texttt{gmp}.
% END C72M8

% BEGIN C72M9
\paragraph{C72M9. The dual abelian scheme and smooth crystalline comparison.}
Take the relative dual $A^\vee$ with its normalized Poincaré bundle
$\mathcal P$ by Artin--Raynaud dual representability (Faltings--Chai,
\emph{Degeneration of abelian varieties}, Ch.~I, \S1). This applies
to an abelian scheme over the given base without a chosen polarization.
Duality of abelian schemes, the theorem of the square
and biduality identify both factors with the other's dual, with the
same relative dimension. To check the universal property in the
given formulation, let $M$ on $A\times_S T$ be additive in the first
variable. Pulling back along two zero sections gives
$e^*M\simeq e^*M\otimes e^*M$, hence $e^*M$ is trivial. Choose
a trivialization. Additivity on every test scheme implies that
the fibre restrictions are algebraically trivial: on each geometric
fibre the translation-difference homomorphism vanishes, and its kernel
on the Picard group is $\operatorname{Pic}^0$. The rigidified bundle
therefore gives a section of the dual via the relative Picard
functor, hence $M\simeq(1,x)^*\mathcal P$. The section $x$ only
depends on the relative class, so it is unique. Conversely the
biextension law and its normalization give the stipulated additivity
and triviality. Biduality applies to the swapped bundle. Cocycles
on trivializing opens represent these line bundles and isomorphisms
in the source's presentation.

The other half of \texttt{dual} concerns every smooth,
quasi-compact separated lift, including those without a group law.
For $X_0=X\times_R R_0$ D72P2 gives the comparison by restriction
of the bounded divided-power thickenings attached to the actual
immersion. Its isomorphisms in degrees one and two are respectively
the two specified cochain pullbacks. This proves both conjuncts.
% END C72M9

% BEGIN C72M10
\paragraph{C72M10. The integral Hodge filtration and its polarization form.}
Write $H=H^1_{\rm dR}(A/R)$, $\omega=e^*\Omega^1_{A/R}$. The
abelian-scheme cohomology theorem (Berthelot--Breen--Messing,
\emph{Théorie de Dieudonné cristalline II}, Proposition~2.5.2)
gives products $\bigwedge^j H\simeq H^j_{\rm dR}$, degeneration of
the Hodge spectral sequence and arbitrary base change. In particular
\[
 0\longrightarrow\omega\longrightarrow H
 \longrightarrow R^1 f_*\mathcal O_A\longrightarrow0
\]
has locally free outer terms of rank $g$. On the affine base the
sequence of modules splits because the quotient is projective.
Addition pulls each degree-one class back to the sum of its two
projections, as is seen by Künneth and restriction to the axes.
Pulling this identity along any pair of maps proves the primitive
condition even when their common domain is singular.
D72P3 identifies the indicated Hodge submodule with $\omega$, and
evaluation of the form cocycle on the two dual-number lifts with
evaluation on the tangent vector. Dualizing the displayed sequence
gives a surjection $H^\vee\to\omega^\vee=\operatorname{Lie}(A)$
with kernel $\operatorname{Ann}(\omega)$. This gives existence and
uniqueness in the two tangent tests. Base change carries $\omega$
to the span of its image. D72P2 gives the required bijection on the
primitive modules under change from a smooth lift to its reduction:
primitivity is respected in both directions by applying the same
comparison to $A\times A$ and testing the addition identity there.

For any line bundle $L$, the logarithmic first Chern class belongs
to the first Hodge step of $H^2$. Indeed it is represented on
overlaps by $d\log g_{ij}$, where $g_{ij}$ are transition units.
On the basis of tests subordinate to a trivialization, the
source takes the divided-power log of the coboundary of lifted
transition units. This computes the class on the whole site by
descent in D72P2. The integral formula
$\log_{\rm PD}(1+a)=\sum_{j\ge1}(-1)^{j-1}(j-1)!\gamma_j(a)$
is finite on each bounded test. In divided-power differentials
$d\gamma_j(a)=\gamma_{j-1}(a)\,da$, so its differential is
$(1+a)^{-1}da$, also in positive characteristic. In the comparison
double complex this carries the log coboundary to the cocycle of
logarithmic differentials of transitions, up to a fixed sign;
changing lifts adds a coboundary. Under the exterior-square identification this first
step is spanned by $\omega\wedge H$. Consequently every Chern form
as defined by the source vanishes on
$\operatorname{Ann}(\omega)\times\operatorname{Ann}(\omega)$.

For perfectness let $\lambda_L:A\to A^\vee$ be the homomorphism
$x\mapsto[t_x^*L\otimes L^{-1}]$. The source's trivial-kernel
test uses points over all affine test schemes, modulo bundles locally
pulled back from their bases. It therefore says that the
scheme-theoretic kernel is zero. The homomorphism is an isomorphism:
on geometric fibres it is a monomorphism of proper smooth connected
groups of equal dimension, hence an isomorphism, and the fibrewise
isomorphism criterion applies to the flat finitely presented groups.
The cross term of the alternating Chern form between the Hodge
annihilator and its quotient is the differential
$d\lambda_L:\operatorname{Lie}(A)\to\operatorname{Lie}(A^\vee)
=R^1 f_*\mathcal O_A$, up to sign and transposition. Indeed translating
$g_{ij}$ by the infinitesimal tangent vector $v$ changes it by
the multiplicative ratio $1+\varepsilon\,v(g_{ij})/g_{ij}$.
This represents both the tangent class of $\lambda_L(v)$ in
$R^1 f_*\mathcal O_A$ and the contraction of the first Hodge
graded part of $c_1(L)$ with $v$. In fact, with $Q=R^1f_*\mathcal O_A$,
the annihilator is $Q^\vee$, its quotient in $H^\vee$ is
$\omega^\vee$, and that contraction maps $\omega^\vee\to Q$.
Thus the cross term and its
transpose are invertible. In a local splitting the diagonal block
on the annihilator is zero, so the whole alternating form is
perfect, including in residue characteristic two. This proves all
five parts of \texttt{dr}.
% END C72M10

% BEGIN C72M11
\paragraph{C72M11. Lifting an abelian scheme and its homomorphisms.}
Use Grothendieck, \emph{Groupes de Barsotti--Tate et cristaux de
Dieudonné}, Ch.~V, \S4, pp.~114--116, the deformation theorem
for abelian schemes. For a thickening $S_0\subset S$ with locally
nilpotent divided powers it gives an equivalence between abelian
schemes on $S$ and pairs $(A_0,W)$, where $W$ is a locally direct
summand of the contravariant crystal $H^1_{\rm cris}(A_0/S)$ lifting
$\omega_{A_0}$. Morphisms are the special-fibre homomorphisms whose
crystalline pullbacks preserve the submodules. Grothendieck defines
this crystal on the nilpotent crystalline site and identifies its value
at a lift with relative de Rham cohomology. This theorem allows an
arbitrary base. The source's uniform vanishing of all products of
divided powers of total degree at least $N$ implies the required
local nilpotence and also $I^N=0$.

On the given affine base the evaluation is finite projective and
its dual is the covariant module used in the source (C72M10 and
D72P2). The annihilator of $W$ is a direct summand in this dual,
reduces to the annihilator of $\omega_{A_0}$ and is carried back
to $W$ by taking annihilators twice. Thus the canonical comparison
is bijective, and the filtration on an existing lift satisfies the
stated admissibility test. Conversely an admissible summand $F$ in
the dual evaluated via the supplied reference lift has exactly the
prescribed reduction: correspondence is evaluation after reduction
on the base-changed degree-one module. Take its annihilator, apply
essential surjectivity of the equivalence, and dualize again. The
lift thereby obtained has exactly the stipulated inverse-image
filtration, and has relative dimension $g$ because the base thickening
has the same underlying points. For a special-fibre homomorphism,
covariant pushforward carrying the source annihilator into the
target annihilator is equivalent by duality to crystalline pullback
preserving the Hodge submodules. Full faithfulness gives the unique
homomorphism over $R$ with the required restriction. These are
\texttt{Canonical}, \texttt{LiftObj} and \texttt{LiftHom} in
\texttt{gm}.
% END C72M11
